\section{Parte 1: Identificación de funcionalidades}

Revisando los módulos que definimos para el SIGC-CUSCO en el laboratorio anterior, identificamos las siguientes funcionalidades principales que debe cubrir el sistema:

\begin{enumerate}[nosep]
    \item Registro e inicio de sesión de usuarios con roles (administrador, docente, participante).
    \item Gestión de instituciones (municipalidades, empresas de turismo).
    \item Creación de capacitaciones (título, fechas, horas, cupos, docente).
    \item Inscripción de participantes con validación de DNI.
    \item Control de asistencia por código QR en cada sesión.
    \item Registro de notas y cálculo de aprobación.
    \item Generación de certificados en PDF con código QR.
    \item Verificación pública de certificados (portal web).
    \item Reportes y exportación de datos a Excel.
    \item Notificaciones por correo (confirmación de inscripción, certificado listo).
    \item Marcado de asistencia offline (sincronización posterior).
    \item Catálogo público de cursos disponibles.
\end{enumerate}
